\section{La trayectoria académica de Wang He (王鹤): del video a los datos sintéticos}

La sección anterior trató sobre la comunidad académica; esta trata sobre la trayectoria personal de Wang He, porque sus proyectos de doctorado anticiparon casi por completo las dos líneas principales que después seguiría la inteligencia corpórea: aprender relaciones de interacción a partir de videos de personas y resolver con datos sintéticos la escasez de anotaciones reales. El primer trabajo consistió en aprender, a partir de videos de personas, interacciones de varios pasos entre personas y objetos, para generar animaciones o accionar una demo sencilla de una taza inteligente. En esencia, debía aprender las relaciones causales entre el estado, la acción, el cambio de estado del objeto y la acción probable siguiente.

Esto ocurrió muy pronto, en 2016. En ese momento no existían los grandes modelos multimodales de hoy ni datos VLA consolidados. Para saber si la taza estaba vacía o llena, si la tapa de la botella estaba abierta o cerrada, si la mano estaba vacía o sostenía un objeto, y dónde estaban los límites de cada segmento de acción, había que dividir el problema en varios modelos, con anotación manual y un sistema complejo. Wang He concluyó más tarde que, aunque aprender un modelo del mundo únicamente a partir de video todavía no es hoy la técnica que más impulsa la inteligencia corpórea, ese proyecto le hizo ver muy pronto que el problema de fondo no es el reconocimiento, sino aprender, a partir de la interacción, cómo cambia el estado del mundo.

\begin{figure}[H]
\centering
\includegraphics[width=0.96\textwidth]{academic-path.png}
\caption{Trayectoria académica de Wang He: de los videos de personas y los modelos del mundo al objetivo de investigación en robots domésticos. Diagrama conceptual propio, elaborado a partir de la conversación de 00:41:15--01:25:08.}
\end{figure}

\begin{knowledgebox}{Lectura de la figura: dos proyectos de doctorado que corresponden a dos líneas principales de la inteligencia corpórea}
La primera línea consiste en comprender, a partir de video, las interacciones causales de varios pasos entre personas y objetos, algo cercano al problema actual del world model; la segunda es la estimación de pose de objetos a nivel de categoría, que usa datos sintéticos o de realidad mixta para suplir la falta de anotaciones reales. Al leer la figura, conviene verlas como dos caras del mismo objetivo: comprender el estado de la interacción y obtener a bajo costo datos con los que se pueda entrenar.
\end{knowledgebox}

\subsection{Del modelado físico a la investigación en IA}

Esta subsección explica por qué Wang He pudo abstraer interacciones complejas en modelos. Su formación de licenciatura y de los primeros años de doctorado se inclinaba hacia la física, los dispositivos y el modelado matemático; estaba acostumbrado a ajustar modelos a partir de datos experimentales y luego usarlos para predecir casos nuevos. Más adelante se dio cuenta de que, en el clean room de semiconductores, la iteración experimental era lenta y el control deficiente, lo cual no correspondía con su ritmo de pensamiento, así que se orientó hacia la IA, la graficación por computadora y la interacción física. Esta experiencia muestra que la inteligencia corpórea no consiste solo en escribir código: también exige comprender los estados de los objetos, la causalidad de las acciones y los cambios físicos.

Lo más difícil del primer proyecto de doctorado no fue una arquitectura de red en particular, sino formulate research problem: qué estados había que extraer, qué acciones modificaban qué estados de los objetos y cómo encadenar esos cambios en un modelo causal temporal. Wang He mencionó que, cuando todavía no conocía la expresión, después popular, Perception-Action Loop, ya le había dibujado a su asesor un diagrama de state-action-change world-state. Ese es precisamente el ciclo de base que más tarde adoptaría la inteligencia corpórea.

\begin{knowledgebox}{Términos clave: modelado de la interacción}
\begin{tabularx}{\textwidth}{p{0.20\textwidth}p{0.34\textwidth}X}
\toprule
Término & Significado & Función en este episodio \\
\midrule
Object State & Estado del objeto, como vacío/lleno, abierto/cerrado, posición/orientación & Determina si la acción siguiente es viable. \\
Action Grammar & Método tradicional que describe con reglas el orden de las acciones & Es interpretable, pero tiene una capacidad limitada de generalización y de aprovechar los datos. \\
World Model & Modelo que predice cómo las acciones cambian el estado del mundo & Uno de los objetivos de largo plazo de la inteligencia corpórea. \\
Perception-Action Loop & Ciclo de percepción, acción y retroalimentación & Convierte el reconocimiento estático en inteligencia de ciclo cerrado. \\
\bottomrule
\end{tabularx}
\end{knowledgebox}

\subsection{Estimación de pose a nivel de categoría y datos sintéticos}

El segundo trabajo está más cerca de la vía actual de datos para la inteligencia corpórea. La 6D pose estimation tradicional suele dirigirse a una instancia conocida: primero se define para el objeto un origen de coordenadas y una orientación estándar, y luego se predicen su rotación y su traslación respecto de ese sistema de coordenadas. El problema que planteó Wang He fue la estimación de pose a nivel de categoría: ante una taza que nunca ha visto, ¿puede el modelo saber que la boca de la taza apunta hacia arriba y que el asa está a la derecha, sin tener que escanear y modelar cada taza ni definir un sistema de coordenadas para cada una?

Los datos reales prácticamente no existían. Las imágenes de internet no tienen anotaciones de pose en seis dimensiones, y un estudiante de doctorado no puede fotografiar ni anotar una cantidad ilimitada de tazas. Por eso recurrió a datos de graficación por computadora y de realidad mixta: fue a IKEA a fotografiar fondos reales de mesas con su profundidad, renderizó tazas digitales sobre esos fondos reales, obtuvo automáticamente las etiquetas de pose y generó cientos de miles de imágenes de entrenamiento. Este método coincide en gran medida con las ideas actuales de datos sintéticos y Sim2Real: el fondo real aporta anclaje en el mundo real, el primer plano virtual aporta etiquetas controlables y, al final, la evaluación se hace en el mundo real.

\begin{importantbox}{Una versión temprana de los datos sintéticos}
La estimación de pose a nivel de categoría muestra que los datos sintéticos no son una palabra de moda surgida apenas hoy. Siempre que la anotación de datos reales no puede llevarse a gran escala, los investigadores recurren a la graficación por computadora, la simulación o la realidad mixta para fabricar muestras de entrenamiento controlables. Los datos sintéticos actuales para robots solo extienden esta idea a más objetos, acciones, escenas e interacciones físicas.
\end{importantbox}

\subsection{Por qué insistir en los robots domésticos}

Esta subsección lleva la trayectoria personal hacia las decisiones estratégicas. Alrededor de 2020, Li Kaifu (李开复) le sugirió a Wang He aplicar sus capacidades en visión tridimensional y estimación de pose a la conducción autónoma, el LiDAR y la estimación de pose de vehículos; Wang He, en cambio, insistió en trabajar en home robot, porque consideraba que el entorno doméstico ofrece el espacio de interacción más amplio y es el más cercano a un «agente inteligente general en el mundo físico». En ese momento era una apuesta muy a contracorriente: se pensaba que los robots domésticos todavía estaban muy lejos, y en China casi no había aliados.

Tras volver a China, Wang He impulsó la línea de inteligencia corpórea en la Universidad de Pekín (北大) y en el BAAI (智源), escribió artículos para presentar la embodied AI y, junto con unos pocos investigadores como Lu Cewu (卢策武), introdujo este concepto en la discusión académica e industrial del país. El punto didáctico aquí no es la leyenda personal, sino la estructura de una elección temprana de dirección: si un investigador cree en la next wave, tiene que soportar un periodo sin consenso, sin aliados y sin interés del capital.

\subsection{Resumen de la sección}

La trayectoria académica de Wang He dejó al descubierto de antemano dos problemas centrales de la inteligencia corpórea: primero, la inteligencia debe comprender cómo las acciones cambian el mundo; segundo, la falta de datos reales debe suplirse con datos sintéticos o de realidad mixta. El objetivo de los robots domésticos, por su parte, lleva estos problemas de investigación hacia aplicaciones más amplias en el mundo físico.
